\documentclass[11pt,a4paper]{article}
\usepackage[T1]{fontenc}
\usepackage[utf8]{inputenc}
\usepackage{lmodern,amsmath,amssymb,amsthm,mathtools,booktabs,longtable,microtype}
\usepackage[margin=27mm]{geometry}
\usepackage{xurl,xcolor}
\usepackage[colorlinks=true,linkcolor=blue!50!black,citecolor=blue!50!black,urlcolor=blue!50!black]{hyperref}
\hypersetup{pdftitle={One Seed for Every Adaptive Policy: Exact Causal Simulation and Sharp Entropy Asymptotics},pdfauthor={Jan Mikulik}}
\newtheorem{theorem}{Theorem}[section]
\newtheorem{lemma}[theorem]{Lemma}
\newtheorem{corollary}[theorem]{Corollary}
\newtheorem{proposition}[theorem]{Proposition}
\theoremstyle{definition}\newtheorem{definition}[theorem]{Definition}
\theoremstyle{remark}\newtheorem{remark}[theorem]{Remark}
\newcommand{\E}{\mathbb E}
\newcommand{\Pp}{\mathbb P}
\newcommand{\Prb}{\mathbb P}
\newcommand{\R}{\mathbb R}
\newcommand{\N}{\mathbb N}
\newcommand{\cS}{\mathcal S}
\newcommand{\At}{\operatorname{At}}
\newcommand{\st}{\preceq_{\mathrm{st}}}
\newcommand{\cstar}{c_\star}
\newcommand{\norm}[1]{\left\|#1\right\|_1}
\newcommand{\TV}{\mathrm{TV}}
\newcommand{\repo}{https://github.com/jendamikulik/finite-Shannon-theorem}
\numberwithin{equation}{section}
\setlength{\emergencystretch}{2em}
\title{One Seed for Every Adaptive Policy:\\Exact Causal Simulation and Sharp Entropy Asymptotics}
\author{Jan Mikul\'ik}
\date{13 October 2026}
\begin{document}
\maketitle
\begin{abstract}
We study the least entropy of a random seed that simulates a controlled
history tree exactly under every deterministic adaptive policy, using a
single decoder consistent at shared histories. For a finite tree, let $L$
be the mean of the pointwise lower envelope of the policy information
CDFs. A causal greedy construction gives $L\le C\le L+c_\star$, where
$c_\star=(1+\log_2 e)/2$. We extend the result to locally finite infinite
trees: $C_\infty=\sup_n C_n$, with attainment whenever this value is finite.
We also identify the attained atomic capacity and construct one countable
causal seed satisfying simultaneous R\'enyi bounds.
For stationary finite kernels with common entropy $h$, the excess is
$\sqrt{2/\pi}(\sigma_{\max}-\sigma_{\min})\sqrt n+o(\sqrt n)$.
When the information variances coincide at $v>0$, we prove the finer law
\[
 C_n=nh+\frac{\kappa_{\max}-\kappa_{\min}}{3v\sqrt{2\pi e}}\ln n+O(1),
\]
where the $\kappa_i$ are centered third information moments. The proof
combines an explicit Gaussian corrector with uniformly bounded integrated
Bellman amplification. No lattice assumption is required. The main
statements are linked to their Lean formalizations in a fixed repository
revision. The numerical additive constant and the nonlinear diffusion
profile are attributed to prior work.
\end{abstract}
\noindent\textbf{Keywords:} exact simulation; shared randomness; adaptive policies;
minimum entropy coupling; information spectrum; Bellman equation; formal verification.
\setcounter{tocdepth}{1}
\tableofcontents
\clearpage

\section{Introduction and main results}\label{sec:intro}
A simulator receives an action, produces a random response, and repeats.
The next action may depend on the entire observed history. We ask how much
randomness must be drawn at the start so that every later response is a
deterministic function of that seed, the history, and the current action.
Exactness is required for every policy, with the same seed law and decoder.

The compatibility requirement is substantive. Coupling the transcript
laws of all policies does not, by itself, ensure that two policies receive
the same response when they have the same history and choose the same
action. We impose this agreement atom by atom by taking seed values to
specify total response strategies. Entropy measures the information in
the initial seed, not the size of a representation of the decoder or its
running time.

\subsection{The finite theorem}
For a finite controlled tree $T$, let $I_\pi$ be the information, in bits,
of its random transcript under policy $\pi$. Write
\[
 F_*(t)=\min_\pi\Pp(I_\pi\le t),\qquad
 L(T)=\int_0^\infty(1-F_*(t))\,dt.
\]
Let $C(T)$ be the minimum entropy of an exact causal seed. A seed may
be finite or countable; these give the same optimum on a finite tree.

\begin{theorem}[Finite causal Shannon theorem]\label{thm:finite}
Every finite controlled tree admits a finite exact causal seed $S_G$ with
\begin{equation}\label{eq:sandwich}
 L(T)\le C(T)\le H(S_G)\le L(T)+c_\star,
 \qquad c_\star=\frac{1+\log_2 e}{2}.
\end{equation}
More precisely, if $Z$ has CDF $F_*$ and $J_G$ is the information of $S_G$,
then
\begin{equation}\label{eq:stoch}
 Z\st J_G\st Z+W,\qquad
 W=A(1+E),\quad A\sim\mathrm{Bernoulli}(1/2),\quad
 E\sim\mathrm{Exp}(\ln2),
\end{equation}
where $Z,A,E$ on the right are independent.
\end{theorem}
The bound is independent of depth and of all finite local alphabet sizes.
The constant $c_\star$ is the ordinary greedy-coupling constant of
Compton et al.~\cite{compton}; the proof below supplies a compatible
strategy construction on a controlled tree.

\subsection{Infinite horizons}
Let $T_n$ be the depth-$n$ truncation of a locally finite tree, and let
$C_\infty$ be the infimum of the entropy of one countable seed exact at
every finite horizon. An empty feasible class has cost $+\infty$.
\begin{theorem}[Projective identity]\label{thm:infinite}
\[
 C_\infty=\sup_n C(T_n),\qquad
 L_\infty\le C_\infty\le L_\infty+c_\star,
 \quad L_\infty=\sup_n L(T_n).
\]
If $C_\infty<\infty$, a countable exact seed attains the infimum.
Consequently a finite-entropy exact seed exists if and only if
$\sup_n L(T_n)<\infty$.
\end{theorem}
Section~\ref{sec:infinite} also treats atomic capacity when the best exact
representation has a nonatomic component. Section~\ref{sec:renyi} gives
one countable causal seed with simultaneous bounds for every positive
R\'enyi order. These are distinct assertions: countable support need not
imply finite Shannon entropy, and a simultaneous approximation need not
attain the Shannon minimum.

\subsection{The stationary hierarchy}
Suppose a fixed finite family $(P_i)$ is available at every history, with
common entropy $h$. Put
\[
 X_i=-\log_2P_i(Y_i)-h,\quad Y_i\sim P_i,\quad
 \sigma_i^2=\E X_i^2,\quad \kappa_i=\E X_i^3.
\]
Let $C_n$ denote the optimal seed entropy for $n$ rounds.
\begin{theorem}[Stationary entropy asymptotics]\label{thm:hierarchy}
Set $a=\max_i\sigma_i$ and $b=\min_i\sigma_i$.
\begin{enumerate}
\item For $a>b\ge0$,
\[
 C_n=nh+\sqrt{2/\pi}(a-b)\sqrt n+o(\sqrt n).
\]
If $b>0$, the remainder is $O(n^{1/4}\ln n)$.
\item If every $\sigma_i^2=v>0$, then
\begin{equation}\label{eq:logmain}
 \boxed{C_n=nh+c\ln n+O(1),\qquad
 c=\frac{\kappa_+-\kappa_-}{3v\sqrt{2\pi e}},}
\end{equation}
where $\kappa_+=\max_i\kappa_i$ and $\kappa_-=\min_i\kappa_i$.
\item If $a=b=0$, then $C_n=nh$ for every $n$.
\end{enumerate}
All asymptotic constants refer to a fixed kernel. In parts 1 and 2,
the same expansion holds for the envelope mean and the greedy seed entropy.
\end{theorem}
The third assertion is elementary: zero information variance means each
law is uniform on its support. Equal entropy makes all support sizes the
same integer $m$. A sequence of $n$ independent uniform labels in
$\{1,\ldots,m\}$, relabeled for each action, is exact with entropy $n\log_2m$.
The fixed-policy entropy lower bound gives equality.

For $v>0$, the first three centered information moments therefore determine
which of the following orders occurs:
\begin{center}
\begin{tabular}{ll}
\toprule
Moment condition (common entropy) & Excess $C_n-nh$ \\
\midrule
Different variances & $\Theta(\sqrt n)$ \\
Common variance, different third moments & $c\ln n+O(1)$, $c>0$ \\
Common variance and third moment & $O(1)$ \\
\bottomrule
\end{tabular}
\end{center}
The last row is a boundedness assertion, not a positive lower bound or a
claim that the bounded remainder converges.

\subsection{Relation to previous work}
Minimum entropy coupling and functional representation provide the
ordinary, marginal version of the problem. Li~\cite{li} treats approximate
couplings for multiple, including infinite, families of distributions.
Shkel and Yadav~\cite{shkel} develop information-spectrum converses.
Compton et al.~\cite{compton} obtain the additive constant $c_\star$ through
the profile method. We use that constant with explicit attribution and
prove the requisite residual-mass estimate for compatible tree strategies.
A later approximation scheme of Compton~\cite{ptas} concerns the ordinary
finite-marginal optimization problem; our construction does not assert a
polynomial running time in a succinct description of the controlled tree.

The controlled diffusion underlying the square-root regime belongs to the
nonlinear central limit theory of Peng~\cite{peng}. The one-sided capacity
profile appears in Peng and Zhou~\cite{pz}; convergence rates under
sublinear expectations are studied by Fang et al.~\cite{fang}. We include
a discrete comparison proof to establish the precise connection with
causal seed entropy and to make the passage to means explicit.
The equal-variance analysis instead resolves the logarithmic displacement
through a third-moment corrector and a bounded amplification lemma.
We do not claim historical novelty for the ordinary coupling bound, the
nonlinear heat profile, or the underlying information-spectrum converse.

\subsection{Conventions and formalization}
All entropies and information variables in this paper are in bits;
$\ln$ is the natural logarithm. In particular $c$ in
\eqref{eq:logmain} is measured per unit of $\ln n$.
We use $0\log 0=0$. Unless stated otherwise, local kernels are positive;
Section~\ref{sec:zero} explains the extension that keeps zero branches.

The exposition is tied to repository revision
\href{https://github.com/jendamikulik/finite-Shannon-theorem/tree/5741ee193bd5ff7f4addd4c09b2d65edfe238c20}{\texttt{5741ee193bd5}}.
Section~\ref{sec:formal} maps the principal statements to Lean declarations
and specifies what the recorded axiom audit establishes. This paper is an
exposition of the selected main results, not a claim that every auxiliary
result in the repository is reproduced here.

\section{Finite controlled trees and exact seeds}\label{sec:finite}
At a history $v$, a finite nonempty set of actions is available. Each
action $a$ has a finite nonempty output set $\mathcal Y(v,a)$ and a
probability kernel $K_v(y\mid a)$. The child history is $vay$. Put
\[
 p(\varnothing)=1,\qquad p(vay)=p(v)K_v(y\mid a).
\]
Terminal histories may occur at different depths. A policy selects an
action at every nonterminal history; a response strategy $d$ selects an
output at every history-action pair. Together they determine one terminal
transcript, denoted $\operatorname{run}(\pi,d)$.

Write $\chi_d(v)$ for the indicator that $d$ agrees with every output
along $v$. For each action, $\sum_y\chi_d(vay)=\chi_d(v)$.
A seed with label weights $(w_j)$ and response strategies $(d_j)$ is exact
when
\begin{equation}\label{eq:exact}
 \sum_j w_j\chi_{d_j}(v)=p(v)\qquad\text{for every history }v.
\end{equation}
This is equivalent to the prescribed transcript law under every policy.
For one direction sum over the terminal continuations of a history under
a policy following its actions; for the other apply \eqref{eq:exact} to
terminal histories compatible with the policy. Choices at histories
incompatible with a strategy still belong to its total decoder.

There are finitely many strategies. Coarsening any countable seed by its
strategy preserves exactness and cannot increase entropy. The exact laws
on strategies form a nonempty compact polytope. Nonemptiness follows by
independently sampling a response at each history-action pair. Continuity
of finite Shannon entropy shows that $C(T)$ is a minimum.

\subsection{An information-spectrum converse}
For every policy, a seed atom of mass $w$ is mapped into a transcript atom
of mass $p\ge w$. Thus seed information $J$ dominates transcript information
under the induced coupling. Consequently
\begin{equation}\label{eq:lowercdf}
 F_J(t)\le F_{I_\pi}(t)\quad\text{for every }\pi,
 \qquad F_J(t)\le F_*(t).
\end{equation}
The minimum $F_*$ is a CDF because there are finitely many policies.
Integration gives $H(S)\ge L(T)$. This converse also holds for ordinary
couplings of the policy transcripts without shared-history consistency.
For any exact seed with finite entropy, the one-dimensional transport
identity gives
\begin{equation}\label{eq:Wgap}
 W_1(\mathcal L(J),\mathcal L(Z))
 =\int_\R|F_J-F_*|=H(S)-L(T).
\end{equation}

\subsection{Bottleneck minimax and extraction}
A residual flow $r$ is nonnegative and satisfies
$\sum_y r(vay)=r(v)$ for every action. Define
\[
 B_r(v)=r(v)\quad(v\text{ terminal}),\qquad
 B_r(v)=\min_a\max_y B_r(vay)\quad(v\text{ nonterminal}).
\]
\begin{lemma}[Tree minimax]\label{lem:minimax}
If $\mathcal L(d)$ and $\mathcal L_\pi$ denote the leaves compatible with
strategy $d$ and policy $\pi$, respectively, then
\[
 \max_d\min_{v\in\mathcal L(d)}r(v)
 =B_r(\varnothing)
 =\min_\pi\max_{v\in\mathcal L_\pi}r(v).
\]
This is the largest $w$ for which $r-w\chi_d$ is a nonnegative flow for
some strategy $d$.
\end{lemma}
\begin{proof}
Induct on the tree. A strategy chooses an output for each action and a
continuation on every selected child. Maximizing its smallest compatible
leaf mass gives $\min_a\max_y B_r(vay)$, since the choices for distinct
actions are independent. A policy chooses one action and continuations
for every output, giving the same recursion for its smallest achievable
largest leaf mass. The leaf case is immediate. Nonnegativity on the
leaves implies nonnegativity on every cylinder by flow conservation.
Subtracting $w\chi_d$ preserves that conservation. This proves the last
assertion as well.
\end{proof}
Starting with $r=p$, extract a maximizing strategy with weight
$w=B_r(\varnothing)$ and subtract $w\chi_d$. Positive root mass gives
positive bottleneck by induction. At least one previously positive leaf
becomes zero. Hence extraction terminates in at most the number of leaves
of the full controlled tree. The extracted weights are nonincreasing,
sum to one, and satisfy \eqref{eq:exact}. This is the seed $S_G$.

\subsection{The refined cutoff and the universal constant}
Fix $0<\delta\le1$ and stop just after all weights greater than $\delta$
have been extracted. The remaining root mass is
$R=\sum_{j:w_j\le\delta}w_j$. By Lemma~\ref{lem:minimax}, some policy has
every residual leaf mass at most $\delta$. On this policy a leaf with
original mass $p>\delta$ must already have lost a weight greater than
$\delta$. Its residual is therefore at most
\[
 a_\delta(p)=\begin{cases}
 p,&p\le\delta,\\
 \min\{\delta,p-\delta\},&p>\delta.
 \end{cases}
\]
Define the nondecreasing function
\[
 \psi(u)=\begin{cases}
 u,&0\le u\le1/2,\\
 1/2,&1/2<u<1,\\
 1,&u\ge1.
 \end{cases}
\]
For every $p>0$, $a_\delta(p)\le p\psi(\delta/p)$: equality holds
outside $\delta<p<2\delta$, and on that interval the inequality is
$p-\delta\le p/2$. Thus
\begin{align}
 \sum_{j:w_j\le\delta}w_j
 &\le\max_\pi\sum_{v\in\mathcal L_\pi}a_\delta(p(v))\notag\\
 &\le\max_\pi\E\psi(\delta2^{I_\pi})
 \le\E\psi(\delta2^Z).\label{eq:cutoff}
\end{align}
The last step uses $I_\pi\st Z$. It is distinct from the preceding
pointwise majorization: $a_\delta(p)/p$ itself is not monotone in information.

For the variable $W$ in \eqref{eq:stoch}, its inclusive survival at $s$ is
$1$ for $s\le0$, $1/2$ for $0<s<1$, and $2^{-s}$ for $s\ge1$.
Setting $\delta=2^{-t}$ in \eqref{eq:cutoff} therefore gives
\[
 \Pp(J_G\ge t)\le\E\psi(2^{Z-t})=\Pp(Z+W\ge t).
\]
Passing to thresholds decreasing to $t$ proves the strict-tail inequality,
including atoms. Together with \eqref{eq:lowercdf}, this proves
\eqref{eq:stoch}. Finally
$\E W=(1+1/\ln2)/2=c_\star$, completing
Theorem~\ref{thm:finite}.

\begin{remark}
The constant is a universal upper bound, not an overhead paid by every
instance. If all actions request one fixed law, up to bijective
relabeling, an independent sequence with that law is an optimal seed and
$C_n=nh$. Nor does the theorem assert that a minimum-entropy ordinary
coupling is compatible; it constructs a separate causal seed.
\end{remark}

\section{Infinite horizons and atomic capacity}\label{sec:infinite}\subsection{Strategy space and capacity}
Let $T$ be a rooted controlled history tree. At every finite history $h$
there are finitely many available actions, each with a finite nonempty
output alphabet and a prescribed probability kernel. We assume the
alphabets are finite at every node, not uniformly bounded. Terminal nodes
can be padded by a forced action and output. Zero-probability outputs may
be deleted; deterministic singleton alphabets are permitted.

Write $p(h)$ for the product of output probabilities along $h$.
For every available action $a$,
\[
 p(\varnothing)=1,\qquad \sum_y p(hay)=p(h).
\]
At each finite depth there are finitely many histories. A deterministic
response strategy selects one output at every history--action pair.
The space $\cS$ of all such strategies is a countable product of finite
discrete spaces, hence compact and metrizable. For a history $h$, let
$E_h\subseteq\cS$ be the set of strategies agreeing with its output
choices. This is a clopen cylinder, and
\[
 E_h=\bigsqcup_y E_{hay}\quad\text{for every action }a.
\]
A Borel probability measure $\mu$ on $\cS$ is an exact causal
representation if $\mu(E_h)=p(h)$ for every finite history $h$.
This is equivalent to exactness for every deterministic adaptive policy
at every finite horizon, and therefore for the corresponding infinite
path laws.

Let $P_{n,\pi}$ be the finite transcript law at horizon $n$ under policy
$\pi$. Put
\begin{align}
 A_K(P)&=\max_{D:\,|D|\le K}P(D),\\
 a_K(T)&=\inf_{n,\pi} A_K(P_{n,\pi}),\qquad
 q(T)=\lim_{K\to\infty}a_K(T).\label{inf:eq:q}
\end{align}
The limit exists by monotonicity and belongs to $[0,1]$.
For a Borel probability on a compact metrizable space define its total
atomic mass by
\[
 \At(\mu)=\sum_{s:\,\mu(\{s\})>0}\mu(\{s\}).
\]
There are at most countably many terms.

Ordinary couplings below are Borel probability laws on the countable
product of the transcript spaces indexed by all pairs $(n,\pi)$, with
the specified coordinate marginals. They need not be compatible with a
single strategy or even consistent between horizons.

\begin{theorem}[Atomic capacity]\label{inf:th:main}
For every such tree,
\[
 \boxed{\quad
 \max_{\mu\ \mathrm{exact\ causal}}\At(\mu)
 =\max_{\Gamma\ \mathrm{ordinary}}\At(\Gamma)
 =q(T).\quad}
\]
Both maxima are attained. In particular, an exact countably supported
causal seed for the entire infinite tree exists if and only if $q(T)=1$.
If $q(T)<1$, every exact causal representation has nonatomic mass at
least $1-q(T)$, and a representation attaining this minimum exists.
\end{theorem}

\subsection{A finite concentration estimate}
We use the bottleneck construction proved in Section~\ref{sec:finite}. After a
finite number of causal greedy extractions, the residual cylinder masses
form a nonnegative controlled flow. If its root mass is $R$, its next
extractable strategy weight $w$ equals the smallest, over policies, of
the largest residual transcript mass. Extraction terminates on each
finite tree and gives an exact finite seed.

\begin{lemma}\label{inf:lem:concentration}
Suppose every policy on a finite controlled tree has $K$ transcripts
carrying total mass at least $a$. The first $M$ greedy atoms, padded with
zeros if extraction has ended, have total mass at least
\[
 a\bigl[1-(1-1/K)^M\bigr].\label{inf:eq:head}
\]
Consequently the same is true of the $M$ largest seed atoms.
\end{lemma}
\begin{proof}
Let $R_m$ be the residual root mass after $m$ steps, with $R_0=1$.
The minimax identity supplies a policy all of whose residual transcript
masses are at most the next weight $w_{m+1}$. On this policy select at
most $K$ original transcripts of mass at least $a$. The residual mass
outside this set is at most $1-a$, so
\[
 R_m-(1-a)\le K w_{m+1}.
\]
Since $R_{m+1}=R_m-w_{m+1}$,
\[
 R_{m+1}-(1-a)\le(1-1/K)(R_m-(1-a)).
\]
Iteration gives $R_M\le1-a+a(1-1/K)^M$. The same bound holds after
termination. Subtract from one. For $K=1$ and $M\ge1$ the displayed
formula has its usual value $a$.
\end{proof}

\subsection{Attainment of atomic capacity}
\begin{proof}[Proof of Theorem~\ref{inf:th:main}]
\emph{Upper bound, even without causality.}
Take any $m$ atoms of an ordinary coupling, of total mass $b$. Their
images in any one transcript coordinate lie in at most $m$ transcripts,
so $b\le A_m(P_{n,\pi})$ for every $(n,\pi)$. Hence $b\le a_m\le q$.
Exhausting all its atoms gives $\At(\Gamma)\le q$. The same argument
applies directly to a causal representation by projecting each atom
through a policy.

\emph{Limits of weights and strategies.}
For each horizon $n$, take a finite greedy causal seed, order its weights
as $w^{(n)}_1\ge w^{(n)}_2\ge\cdots\ge0$, and pad by zeros.
Extend each finite strategy arbitrarily to an element $s^{(n)}_i$ of
$\cS$. By compactness and diagonal selection there is one subsequence
$n_j\to\infty$ such that, for every $i$,
\[
 w_i^{(n_j)}\to w_i,\qquad s_i^{(n_j)}\to s_i.
\]
Set $\nu=\sum_iw_i\delta_{s_i}$ and $A=\sum_iw_i\le1$; repeated
limit strategies are simply merged as measures. For fixed $K,M$,
Lemma~\ref{inf:lem:concentration}, with $a=a_K(T)$, gives
\[
 \sum_{i=1}^M w_i
 \ge a_K(T)\bigl[1-(1-1/K)^M\bigr].
\]
First let $M\to\infty$ and then $K\to\infty$. It follows that
$A\ge q$.

\emph{No invalid passage of infinitely many atoms through a limit.}
For fixed $h$ and $M$, convergence of strategies makes membership in
the clopen set $E_h$ eventually constant for each of the first $M$
indices. For large $j$, the finite seed is exact at $h$. Thus
\[
 \sum_{i=1}^M w_i\mathbf1_{E_h}(s_i)
 =\lim_j\sum_{i=1}^M w_i^{(n_j)}\mathbf1_{E_h}(s_i^{(n_j)})
 \le p(h).
\]
Monotone convergence in $M$ proves $\nu(E_h)\le p(h)$. Define
\[
 r(h)=p(h)-\nu(E_h)\ge0.
\]
Since the $E_{hay}$ partition $E_h$ for each action,
$\sum_y r(hay)=r(h)$ and $r(\varnothing)=1-A$.

\emph{Realization of the missing mass.}
If $A<1$, assign, independently at every history--action pair, an output
with probabilities $r(hay)/r(h)$ whenever $r(h)>0$, and any distribution
when $r(h)=0$. The countable product law $\lambda$ of these response
choices satisfies
\[
 \lambda(E_h)=r(h)/(1-A).
\]
Indeed this follows by induction along $h$ from the displayed ratios;
after a zero residual prefix both sides vanish. Hence
$\mu=\nu+(1-A)\lambda$ is exact. If $A=1$, simply set $\mu=\nu$.
As $\nu$ is atomic of mass $A$, $\At(\mu)\ge A\ge q$.
The upper bound already proved gives $\At(\mu)=A=q$. In particular,
the residual law in this construction has no atoms when $q<1$.

Finally map $s\in\cS$ to its full tuple of finite-policy transcripts.
This is a measurable map into the countable product of finite spaces.
The image of $\mu$ is an ordinary coupling and has atomic mass at least
$q$, since images of atoms remain atoms. The upper bound makes its
atomic mass exactly $q$. This proves attainment on both sides.
\end{proof}

\subsection{Information spectrum and entropy attainment}
All logarithms in this section have base two. Define
\[
 F_n(t)=\min_\pi P_{n,\pi}\{-\log P_{n,\pi}(X)\le t\},\quad
 F_\infty(t)=\inf_n F_n(t),\quad t\ge0.
\]
Refining a transcript cannot reduce its information, and policies extend
to longer horizons. Therefore $F_n(t)$ decreases with $n$.
For an individual finite law $P$, writing its information CDF as $F_P$,
\[
 F_P(t)\le A_{\lfloor2^t\rfloor}(P),\qquad
 A_K(P)\le F_P(t)+K2^{-t}.
\]
Taking infima over all finite-policy laws proves
\[
 \boxed{q(T)=\lim_{t\to\infty}F_\infty(t).}
\]
Thus atomic capacity is precisely the mass that does not escape to
infinite information in the uniform policy envelope. This statement
does not require treating a possibly defective $F_\infty$ as a
probability CDF.

Let $C_n$ be the minimum entropy of an exact seed for the depth-$n$
tree. Let $C_\infty$ be the infimum entropy among countably supported
exact seeds for the full tree, with value $+\infty$ if there is no such
seed. Set
\[
 L_n=\int_0^\infty(1-F_n(t))\,dt,\quad
 L_\infty=\sup_nL_n,\quad c_*=(1+\log e)/2.
\]

\begin{theorem}[No extra entropy at the projective limit]
In the extended real numbers,
\[
 \boxed{C_\infty=\sup_n C_n.}
\]
If the right-hand side is finite, the infimum on the left is attained.
Consequently the finite Shannon theorem implies
\[
 L_\infty\le C_\infty\le L_\infty+c_*.
\]
In particular, $L_\infty<\infty$ is equivalent to the existence of one
finite-entropy exact causal seed for the entire infinite tree.
\end{theorem}
\begin{proof}
Restriction of a seed proves $C_n\le C_{n+1}\le C_\infty$. If
$B=\sup_n C_n=+\infty$, this already proves the equality.
Assume $B<\infty$ and choose a minimizing finite seed at each horizon.
Such minimizers exist: one can minimize entropy over the compact
finite-dimensional polytope of laws on finite strategies, merging
identical strategies if necessary. Order and pad weights as above.
For $i>K$, $w_i^{(n)}\le1/(K+1)$, whence
\[
 \sum_{i>K}w_i^{(n)}\le \frac{B}{\log(K+1)}.\tag{*}
\]
Diagonal compactness gives limiting weights and strategies. The uniform
tail bound (*) forces $\sum_iw_i=1$. The cylinder argument in the
previous proof gives $\nu(E_h)\le p(h)$. Since both are probability
flows, equality follows at each depth by summing over the leaves of
any policy following $h$, or by passing the uniformly small tails
directly through the limit.

The countable label seed $i$ with masses $w_i$ realizes these strategies.
With $0\log0=0$, continuity of each summand and Fatou's lemma give
\[
 H(w)=\sum_i-w_i\log w_i
 \le\liminf_j H(w^{(n_j)})=B.
\]
Together with restriction this yields $C_\infty=B$, with attainment.
Finally the finite theorem gives $L_n\le C_n\le L_n+c_*$;
take suprema. No continuity of Shannon entropy on an unrestricted
countable simplex is being assumed.
\end{proof}

\paragraph{Three regimes.}
If $q<1$, every exact representation has an unavoidable nonatomic
component. If $q=1$ but $L_\infty=\infty$, a countable exact seed exists
but all such seeds have infinite entropy. If $L_\infty<\infty$, an
optimal finite-entropy countable seed exists. Countable support is not
the same as finite support.

\subsection{Pointwise versus uniform simulability}
Consider a deterministic waiting chain. At waiting stage $m\ge1$ the
controller can either wait for the next stage or trigger an experiment.
Triggering at stage $m$ produces $m$ independent fair bits and then
stops; there are no choices during the experiment. A policy that never
triggers has a deterministic infinite transcript. The tree is locally
finite, with at most two actions and at most two outputs at every node.

Every individual deterministic policy therefore has a finite-support
full-path law: either one path, or $2^m$ paths for some finite $m$.
Nevertheless, for each fixed $K$, choosing arbitrarily large $m$ yields
\[
 a_K\le K2^{-m},\qquad a_K=0,\qquad q=0.
\]
Hence \emph{every exact common representation is atomless}, even
without a causality requirement. An elementary direct proof is that
an atom of mass $w>0$ would have to map into a transcript of mass
$2^{-m}<w$ for a sufficiently large triggered experiment.
This example illustrates why separate finite-support simulations do
not imply one discrete seed. It is not claimed as a new counterexample
in the general theory of ordinary coupling.

For any prescribed $q_0\in(0,1)$, first choose a deterministic absorbing
branch with probability $q_0$ and the waiting construction with
probability $1-q_0$. Then $a_K=q_0$ for every $K\ge1$: the absorbing
transcript gives the lower bound, and triggering an arbitrarily long
experiment gives the upper bound. Thus the full interval of atomic
capacities is realized even with binary branching.


\section{One seed for all positive R\'enyi orders}\label{sec:renyi}
Fix a rooted controlled history tree with finite nonempty action and output
sets at every history and strictly positive transition probabilities.
Finite horizons therefore have finitely many histories and policies.
Terminal histories may be padded by deterministic steps.

A \emph{causal seed} is a countable probability law on total response
strategies: a strategy specifies an output at every history--action pair.
It is exact if it realizes the prescribed law under every deterministic
adaptive policy at every finite horizon. An \emph{ordinary source} is a
countable random variable with a deterministic readout for each such
policy and horizon; agreement between the readouts at shared histories is
not required. Assume at least one ordinary exact source exists.

All logarithms and entropies are in bits. Write
\[
 O_\alpha=\inf_{U\text{ ordinary exact}}H_\alpha(U),\qquad
 C_\alpha=\inf_{U\text{ causal exact}}H_\alpha(U),
\]
where $H_\alpha(p)=(1-\alpha)^{-1}\log_2\sum_i p_i^\alpha$ for
$\alpha\ne1$, $H_1$ is Shannon entropy, and
$H_\infty(p)=-\log_2\sup_i p_i$. Infinite values are allowed.

\begin{theorem}\label{ren:thm:main}
There exists a single exact countable causal seed $S$ such that, simultaneously
for every $0<\alpha<\infty$,
\[
 \boxed{O_\alpha\le C_\alpha\le H_\alpha(S)\le O_\alpha+d_\alpha,}
\]
where
\[
 d_\alpha=\frac{\log_2\!\left(\frac12+\frac{2^{-\alpha}}{\alpha}\right)}
 {1-\alpha}\quad(\alpha\ne1),\qquad
 d_1=\frac{1+\log_2 e}{2}.
\]
Moreover $H_\infty(S)=C_\infty=O_\infty$.
Here $\infty$ denotes the R\'enyi order, not the time horizon.
\end{theorem}

\subsection{A distributional compactness argument}
For the horizon-$n$ transcript law $P_{n,\pi}$, put
\[
 F_n(t)=\min_\pi P_{n,\pi}\{-\log_2P_{n,\pi}(X)\le t\},
 \qquad F(t)=\inf_nF_n(t).
\]
Zero-probability transcripts carry no mass in this expression.
Extending a policy and refining its transcript shows $F_{n+1}\le F_n$.
Each $F_n$ is a CDF; $F$ is nondecreasing and right-continuous. For the latter,
use $F(t)=\inf_n\inf_{u>t}F_n(u)=\inf_{u>t}F(u)$.

Let $I_U=-\log_2 p_U(U)$ for any countable source.
Deterministic readouts merge atoms, so exactness gives
\begin{equation}\label{ren:eq:lower}
 F_{I_U}(t)\le F_n(t)\quad\hbox{for all }n,t.
\end{equation}
Consequently $F(t)\to1$ as $t\to\infty$ whenever an ordinary countable source
exists. Thus $F$ is the CDF of a finite nonnegative random variable $Z$.
Let $Z_n$ have CDF $F_n$; then $Z_n\st Z$ and $Z_n\Rightarrow Z$.

\paragraph{Finite input.}
The finite theorem of Section~\ref{sec:finite} gives a finite exact causal seed
at each horizon, with information $J_n$, for which
\begin{equation}\label{ren:eq:finite}
 Z_n\st J_n\st Z_n+W,
 \qquad
 W=A(1+E),\quad A\sim{\rm Bernoulli}(1/2),\quad
 E\sim{\rm Exp}(\ln2).
\end{equation}
Here $A,E,Z_n$ are independent. The lower comparison is the deterministic-readout inequality.

\begin{lemma}\label{ren:lem:limit}
There exists an exact countable causal seed $S$ with information $J$ such that
\[
 \boxed{Z\st J\st Z+W,}
\]
with $Z$ and $W$ independent on the right-hand side.
\end{lemma}
\begin{proof}
List the finite greedy atom weights in decreasing order, pad by zeros, and
write them as $p_{n,i}$, $i\ge1$. Extend each finite response strategy
arbitrarily to the full infinite tree. As $Z_n\st Z$, \eqref{ren:eq:finite} gives
$J_n\st Z+W$. For $t\ge0$, at most $\lfloor2^t\rfloor$ atoms have weight
at least $2^{-t}$. Hence
\begin{equation}\label{ren:eq:tight}
 \sum_{i>\lfloor2^t\rfloor}p_{n,i}
 \le \Pp(J_n>t)\le\Pp(Z+W>t).
\end{equation}
This is a uniform tail estimate for the \emph{ranked weights}, with no
entropy assumption.

The full strategy space is a countable product of finite discrete sets.
Diagonal extraction gives a subsequence $n_j\to\infty$ such that, for every
$i$, $p_{n_j,i}\to p_i$ and the associated strategies converge coordinatewise
to $s_i$. Finite partial sums have limits. By \eqref{ren:eq:tight}, their tails
are uniformly arbitrarily small, so $\sum_i p_i=1$ and
$\sum_i|p_{n_j,i}-p_i|\to0$.

Give label $i$ weight $p_i$ and response strategy $s_i$, discarding zero
weights. Fix a finite horizon, policy and transcript. The indicator that
strategy $s_{n_j,i}$ generates that transcript eventually equals its limit
for each fixed $i$, because it depends on finitely many coordinates.
Splitting the weighted sum at a finite rank, and then using the uniform
tail estimate, proves exactness of the limit. This works for every finite
policy law.

For completeness, $J_{n_j}\Rightarrow J$ follows directly from the same
argument. For bounded continuous $f$, write
\[
 \E f(J_{n_j})=\sum_i p_{n_j,i}f(-\log_2p_{n_j,i}).
\]
A fixed positive limiting weight gives convergence of its term; a zero
limiting weight gives a term tending to zero in absolute value. Uniformly
small rank tails control the remaining sum. Thus the expectations converge
to $\sum_{i:p_i>0}p_i f(-\log_2p_i)$.

Finally $Z_{n_j}\Rightarrow Z$ and $Z_{n_j}+W\Rightarrow Z+W$.
First-order stochastic order is closed under weak convergence: test against
bounded continuous increasing functions. Passing to the limit in
\eqref{ren:eq:finite} proves the lemma.
\end{proof}

\subsection{Entropy comparisons}
\begin{proof}[Proof of Theorem \ref{ren:thm:main}]
Take the single seed of Lemma \ref{ren:lem:limit}. For $\alpha\ne1$, define
\[
 R_\alpha(Z)=\frac{1}{1-\alpha}\log_2\E 2^{(1-\alpha)Z}.
\]
The stochastic sandwich and independence give
\begin{equation}\label{ren:eq:renyi}
 R_\alpha(Z)\le H_\alpha(S)\le
 R_\alpha(Z)+\frac{\log_2\E2^{(1-\alpha)W}}{1-\alpha}.
\end{equation}
For $0<\alpha<1$ these follow by monotonicity of the increasing exponential;
infinite moments are permitted. For $\alpha>1$ the exponential is decreasing,
and division by $1-\alpha<0$ reverses the inequality once more. Its
expectation is strictly positive because $Z$ is finite almost surely.

Since $1-\alpha<1$, the exponential moment of $W$ is finite, and
\[
 \E2^{(1-\alpha)W}
 =\frac12+\frac12\,2^{1-\alpha}
   \int_0^\infty(\ln2)\,2^{-\alpha x}\,dx
 =\frac12+\frac{2^{-\alpha}}{\alpha}.
\]
For $\alpha=1$, use the nonnegative means directly:
\[
 \E Z\le H_1(S)\le \E Z+\E W,\qquad
 \E W=\frac12\left(1+\frac1{\ln2}\right)=d_1.
\]
This also agrees with the continuous extension of $d_\alpha$ at $1$.

For any ordinary exact source $U$, \eqref{ren:eq:lower} says $Z\st I_U$.
Therefore $R_\alpha(Z)\le H_\alpha(U)$, and $\E Z\le H_1(U)$.
Taking the infimum over such $U$ in \eqref{ren:eq:renyi} yields
$H_\alpha(S)\le O_\alpha+d_\alpha$. The other inequalities follow from
inclusion of feasible classes and feasibility of $S$.

For the endpoint, fix any ordinary exact $U$. The preceding bound gives
$H_\alpha(S)\le H_\alpha(U)+d_\alpha$ for every $\alpha>1$.
For a countable probability law, $H_\alpha\to H_\infty$ as
$\alpha\to\infty$, whereas $d_\alpha\to0$. Consequently
$H_\infty(S)\le H_\infty(U)$ for every $U$. Taking the infimum and using
$O_\infty\le C_\infty\le H_\infty(S)$ proves equality.
No interchange of an infimum with a limit is needed.
\end{proof}

\begin{corollary}[Simultaneous finiteness]
For the same seed $S$, and for each $0<\alpha\le1$,
\[
 H_\alpha(S)<\infty
 \iff
 \text{some ordinary exact countable source has finite }H_\alpha.
\]
In particular the set of orders with finite R\'enyi entropy is realized
by one causal source, even when its Shannon entropy is infinite.
\end{corollary}
\begin{proof}
The forward implication uses $S$ itself as an ordinary source. The reverse
implication follows from the theorem and the finite constant $d_\alpha$.
\end{proof}

\paragraph{Examples of the bounds.}
The same $S$ satisfies
\[
 H_1(S)\le O_1+1.2213475204\ldots,\qquad
 H_2(S)\le O_2+\log_2(8/5),\qquad H_\infty(S)=O_\infty.
\]
These are simultaneous approximation statements, not exact simultaneous
minimization at finite orders. The construction does not claim that $S$
also attains the Shannon optimum.



\section{Keeping zero-probability branches}\label{sec:zero}
The simulation model can retain every branch of a nonnegative normalized
kernel. For information and moment expressions, only positive atoms are
sampled. Let $T^+$ be the subtree of histories whose path probability is
positive, with zero entries omitted from each local output alphabet.
Normalization makes every reduced alphabet nonempty.

Restriction gives a map $R$ from full strategies to reduced strategies:
at a positive history, keep a positive output, and replace any zero output
by a fixed positive default. Extension gives a map $E$ in the other
direction by choosing defaults at all omitted histories. Both maps are
measurable, and deterministic when applied to seed labels.

For an exact full law, the event that a positive history is reached and
a zero-probability response is selected has probability zero. There are
only countably many history-action-output triples. Outside their union,
restriction therefore preserves every positive-history cylinder.
Induction on the path length shows that $R$ sends every exact full seed
to an exact reduced seed. Conversely $E$ sends any exact reduced seed
to an exact full seed: every omitted path has a first zero transition,
which is never selected along a reached positive history.

Using the same label weights preserves seed entropy in either direction;
merging labels after a deterministic map can only decrease a positive-order
R\'enyi entropy. Hence the full and reduced entropy infima agree, for
finite and infinite horizons. The policy laws have the same positive
transcript masses, so their information envelopes also agree. Applying
the preceding theorems on $T^+$ and extending the resulting strategies
proves their versions on the full tree. Zero branches are retained in the
final decoder, although the proof uses the support reduction.

\section{Stationary kernels and the square-root regime}\label{sec:sqrt}
Use the centered variables $X_i$ of Section~\ref{sec:intro}; put
$a=\max_i\sqrt{\E X_i^2}$, $b=\min_i\sqrt{\E X_i^2}$, and
$M=\max_i\|X_i\|_\infty$. The proof applies to arbitrary finitely supported
centered increments, even if they are not information variables.
The conclusions to prove are
\begin{align}
 C_n&=nh+\sqrt{2/\pi}(a-b)\sqrt n+o(\sqrt n),\label{sq:asymptotic}\\
 C_n&=nh+\sqrt{2/\pi}(a-b)\sqrt n+O(n^{1/4}\ln n)\quad(b>0).
 \label{sq:rate}
\end{align}
\subsection{The discrete Bellman envelope}
For a policy $\pi$, let $S_n^\pi=\sum_{k=1}^n X_{A_k}$, with the fresh
increment drawn after the predictable choice $A_k$. Set
\[
 G_n(x)=\inf_\pi\Prb(S_n^\pi\le x),\qquad
 (Tf)(x)=\min_{j\in\mathcal A}\E f(x-X_j).
\]
Condition on the first action and outcome. Continuation policies can be
chosen separately for each outcome; finite horizon and finite alphabets
ensure the minima are attained. Thus
\begin{equation}\label{sq:bellman}
 G_0(x)=\mathbf1_{\{x\ge0\}},\qquad G_{n+1}=TG_n.
\end{equation}
The operator is monotone, commutes with spatial translations and addition
of constants, and satisfies $\|Tf-Tg\|_\infty\le\|f-g\|_\infty$.
Moreover $G_n$ is a CDF: it is the minimum of finitely many CDFs with bounded
support. By the chain rule for transcript probabilities,
\begin{equation}\label{sq:meanidentity}
 \frac{L_n-nh}{\sqrt n}
 =\int_0^\infty[1-G_n(x\sqrt n)]\,dx
  -\int_0^\infty G_n(-x\sqrt n)\,dx.
\end{equation}
There is no exchange of an infimum with an integral in this identity.
The envelope is formed first, and then its mean is integrated.

\subsection{A quantitative comparison for positive least variance}
Assume $0<b\le a$. Write $\phi(z)=(2\pi)^{-1/2}e^{-z^2/2}$ and let $\Phi$
be its CDF. Define
\begin{equation}\label{sq:profile}
 F_{a,b}(z)=
 \begin{cases}
 \dfrac{2b}{a+b}\Phi(z/b),&z\le0,\\[3pt]
 1-\dfrac{2a}{a+b}\Phi(-z/a),&z\ge0.
 \end{cases}
 \qquad u(t,x)=F_{a,b}(x/\sqrt t),\quad t>0.
\end{equation}
Both values at zero are $b/(a+b)$. On either half-line, with $s=b$ on the
negative side and $s=a$ on the positive side,
\[
 F'(z)=\frac{2}{a+b}\phi(z/s),\qquad
 F''(z)=-\frac{2z}{(a+b)s^2}\phi(z/s).
\]
Consequently $F\in C^2(\R)$, $F''(0)=0$, and
\begin{equation}\label{sq:pde}
 u_t=\min_j\frac{\sigma_j^2}{2}u_{xx}.
\end{equation}
Indeed, $u_{xx}\ge0$ for $x<0$, selecting $b^2$, and $u_{xx}\le0$ for $x>0$,
selecting $a^2$; both sides vanish at zero. No classical third derivative at
zero is required.

Fix $M=\max_j\|X_j\|_\infty$ and $\phi_0=(2\pi)^{-1/2}$. Define
\begin{align*}
 A&=\frac{M^3\phi_0}{3(a+b)b^2},&
 B&=\frac{a\phi_0}{2(a+b)}
       \bigl(3e^{-1/2}+3\sqrt3\,e^{-3/2}\bigr),\\
 \ell&=\frac{2\phi_0}{a+b},&
 D&=\frac{a\phi_0 e^{-1/2}}{a+b}.
\end{align*}

\begin{lemma}[One-step defect]\label{sq:local}
For every $t>0$,
\begin{equation}\label{sq:defect}
 \|Tu(t,\cdot)-u(t+1,\cdot)\|_\infty
 \le A t^{-3/2}+\frac B2 t^{-2}.
\end{equation}
\end{lemma}
\begin{proof}
On each open half-line,
\[
 F'''(z)=\frac{2}{(a+b)s^2}\bigl((z/s)^2-1\bigr)\phi(z/s).
\]
Since $\sup_y|(y^2-1)\phi(y)|=\phi_0$ and $F''$ is continuous at zero,
$F''$ is globally Lipschitz with constant $2\phi_0/((a+b)b^2)$.
Taylor's integral remainder therefore gives, also across zero,
\[
 \left|\E u(t,x-X_j)-u(t,x)-\frac{\sigma_j^2}{2}u_{xx}(t,x)\right|
 \le A t^{-3/2}.
\]
The linear term vanishes because $\E X_j=0$. Taking the minimum changes
neither this bound nor its uniformity in $x$. Equation \eqref{sq:pde} identifies
the quadratic part with $u_t$.

For $z=x/\sqrt t$,
\[
 u_{tt}(t,x)=\frac{3zF'(z)+z^2F''(z)}{4t^2}.
\]
The elementary maxima $\sup_y|y|\phi(y)=\phi_0e^{-1/2}$ and
$\sup_y|y|^3\phi(y)=3\sqrt3\phi_0e^{-3/2}$ give
$\|u_{tt}(t,\cdot)\|_\infty\le B/t^2$. Integrating this bound over a
unit time step supplies the second remainder $B/(2t^2)$.
\end{proof}

\begin{lemma}[Explicit CDF comparison]\label{sq:cdf}
For $n\ge1$, $m\ge1$ and $R\ge0$,
\begin{align}\label{sq:cdfbound}
 \sup_z|G_n(z\sqrt n)-F_{a,b}(z)|\le{}&2\Phi(-R)
   +\frac{3A}{\sqrt m}+\frac Bm\\
 &+\ell aR\sqrt{\frac m{n+m}}
   +D\ln(1+m/n).\nonumber
\end{align}
In particular the left side is $O(n^{-1/4}\sqrt{\log n})$ for $n\ge2$.
\end{lemma}
\begin{proof}
Let $d=aR\sqrt m$ and $\delta=2\Phi(-R)$. Directly from \eqref{sq:profile},
\[
 u(m,x-d)-\delta\le\mathbf1_{\{x\ge0\}}
                       \le u(m,x+d)+\delta.
\]
Apply $T^n$, preserving order and the translations. By the contraction
property and Lemma~\ref{sq:local}, the accumulated defect is at most
\[
 E_m:=\sum_{k=0}^{n-1}\left(\frac A{(m+k)^{3/2}}
                         +\frac B{2(m+k)^2}\right)
 \le\frac{3A}{\sqrt m}+\frac Bm.
\]
Thus $u(n+m,x-d)-\delta-E_m\le G_n(x)
\le u(n+m,x+d)+\delta+E_m$.
The spatial translation costs at most $\ell d/\sqrt{n+m}$.
Also $\sup_x|u_t(t,x)|\le D/t$, so changing time from $n+m$ to $n$
costs at most $D\ln(1+m/n)$. This proves \eqref{sq:cdfbound}.
Choose $m=\sqrt n$, $R=\sqrt{\ln n}$, and use
$2\Phi(-R)\le e^{-R^2/2}$ to obtain the stated order.
\end{proof}

\subsection{The mean: uniform tails, an explicit error bound, and the coefficient}
Under every policy, the increments are martingale differences in $[-M,M]$.
Conditional Hoeffding's bound and iteration give
$\E e^{\theta S_n^\pi}\le e^{n\theta^2M^2/2}$.
Markov's inequality, optimized at $\theta=x/(M^2\sqrt n)$ and applied also
to $-S_n^\pi$, yields, for $x>0$,
\begin{equation}\label{sq:tails}
 1-G_n(x\sqrt n)\le e^{-x^2/(2M^2)},\qquad
 G_n(-x\sqrt n)\le e^{-x^2/(2M^2)}.
\end{equation}
These bounds are uniform in both $n$ and the adaptive policy. The same
bounds hold for $1-F_{a,b}(x)$ and $F_{a,b}(-x)$ since $a\le M$.
Dominated convergence in \eqref{sq:meanidentity} is now justified.
Moreover, if $e_n$ denotes the right side of \eqref{sq:cdfbound}, then, for every $r>0$,
\begin{equation}\label{sq:meanbound}
 \left|\frac{L_n-nh}{\sqrt n}-\mu_{a,b}\right|
 \le 2r e_n+\frac{4M^2}{r}e^{-r^2/(2M^2)},
\end{equation}
where $\mu_{a,b}$ is the mean of $F_{a,b}$. Split the CDF-distance integral
at $\pm r$ to obtain this inequality, using
$\int_r^\infty e^{-x^2/(2M^2)}dx\le(M^2/r)e^{-r^2/(2M^2)}$.
With $r=M\sqrt{2\ln n}$, Lemma~\ref{sq:cdf} makes \eqref{sq:meanbound}
$O(n^{-1/4}\ln n)$.

Finally, $\int_0^\infty\Phi(-x/s)\,dx=s/\sqrt{2\pi}$ gives
\begin{equation}\label{sq:coefficient}
 \mu_{a,b}
 =\frac{2a^2-2b^2}{(a+b)\sqrt{2\pi}}
 =\sqrt{\frac2\pi}(a-b).
\end{equation}
Equations \eqref{sq:meanbound}, \eqref{sq:coefficient} and \eqref{eq:sandwich}
prove both claims of Theorem~\ref{thm:hierarchy}(1) when $b>0$. In particular,
\eqref{sq:cdfbound} and \eqref{sq:meanbound}, plus $\cstar/\sqrt n$, are an
explicit finite-$n$ bound for the normalized optimal entropy error.

\subsection{Zero least variance: a controlled perturbation argument}
Assume $b=0<a$. For $\varepsilon>0$, replace every increment by
$X_j+\varepsilon\xi$, where each fresh $\xi$ is an independent symmetric
sign. The controller is allowed to observe the pair of outcomes. The Bellman
state is still the partial sum: the conditional future laws depend only on
the next action, so observing the pair rather than its sum changes no value.
Let $G_n^\varepsilon$ be this envelope. Its extreme standard deviations are
$a_\varepsilon=\sqrt{a^2+\varepsilon^2}$ and $b_\varepsilon=\varepsilon$;
the previous proof applies to this bounded, centered family.

For every $\eta>0$ and $x\in\R$,
\begin{equation}\label{sq:noisecompare}
 G_n^\varepsilon((x-\eta)\sqrt n)-\varepsilon^2/\eta^2
 \le G_n(x\sqrt n)
 \le G_n^\varepsilon((x+\eta)\sqrt n)+\varepsilon^2/\eta^2.
\end{equation}
To verify this, couple the original and perturbed sums using the same
actions and original outcomes. Their difference is
$\varepsilon\sum_{k=1}^n\xi_k$, whose tail beyond $\eta\sqrt n$ is at most
$\varepsilon^2/\eta^2$. An original policy can be used in the perturbed
system by ignoring the signs. Conversely, a perturbed policy can be run in
the original system using independent auxiliary signs. Such randomization
does not improve an infimum of event probabilities: conditional on its
entire independent random tape it is a deterministic causal policy, so its
event probability is an average of probabilities of admissible policies.
These two comparisons, with shifted thresholds, prove \eqref{sq:noisecompare}.

First let $n\to\infty$ for fixed $\varepsilon,\eta$, then set
$\eta=\sqrt\varepsilon$ and let $\varepsilon\downarrow0$. The explicit
profiles tend to the continuous CDF
\[
 F_{a,0}(x)=\begin{cases}0,&x\le0,\\2\Phi(x/a)-1,&x>0.\end{cases}
\]
Equation \eqref{sq:noisecompare} proves $G_n(x\sqrt n)\to F_{a,0}(x)$ at
every $x$. The original-family bound \eqref{sq:tails} still holds, so dominated
convergence gives mean $a\sqrt{2/\pi}$. This proves \eqref{sq:asymptotic}
for $b=0<a$, without asserting the remainder \eqref{sq:rate} in that case.
If $a=b=0$, every increment is zero almost surely, $G_n=G_0$, and
\eqref{eq:sandwich} immediately gives \eqref{sq:asymptotic}. This finishes the proof.

\subsection{Universality of near-optimal seed information}
\begin{corollary}[Universality of asymptotically optimal seeds]
Suppose $a>0$. Let $J_n=-\log_2 w_n(S_n)$ be the self-information of an
exact countable seed with finite entropy, and suppose $\E J_n-C_n=o(\sqrt n)$.
Then the law of $(J_n-nh)/\sqrt n$ converges in $W_1$ to the law with CDF
$F_{a,b}$ (with the half-normal convention when $b=0$).
\end{corollary}
\begin{proof}
Every seed atom is no heavier than any transcript to which it is mapped.
Thus its self-information stochastically dominates each policy transcript,
and its CDF is at most their lower envelope. If $Z_n$ has that envelope CDF,
then $Z_n\preceq_{\rm st}J_n$, and
$W_1(\mathcal L(J_n),\mathcal L(Z_n))=\E J_n-L_n$.
By \eqref{eq:sandwich} this is $o(\sqrt n)$. The CDF convergence and uniform
tail bound above give $W_1$ convergence of $(Z_n-nh)/\sqrt n$ to $F_{a,b}$.
The triangle inequality proves the claim.
\end{proof}

The profile is not the limiting transcript law of one policy when $a>b$:
every centered transcript has mean zero and the uniform tail bound forces
any weak limit to have mean zero. The envelope profile has strictly positive
mean. Different thresholds may need different minimizing policies. Treating
the whole envelope as one policy law would invalidate the argument.


\section{The logarithmic regime with bounded remainder}\label{sec:log}
Let $X_1,\ldots,X_d$ have finitely supported laws satisfying
\[
 \E X_i=0,\qquad \E X_i^2=v>0.
\]
Set $\kappa_i=\E X_i^3$, $\kappa_-=\min_i\kappa_i$,
$\kappa_+=\max_i\kappa_i$, and
\[
 A_i f(x)=\E f(x-X_i),\qquad
 Tf=\min_i A_i f,\qquad Uf=\max_i A_i f.
\]
Both operators preserve order and commute with translations. The
operator $U$ is positively homogeneous and subadditive, and
\begin{equation}\label{log:eq:lip}
 |Tf-Tg|\le U|f-g|,\qquad |Uf-Ug|\le U|f-g|.
\end{equation}
Let $H=\mathbf1_{[0,\infty)}$, $G_n=T^nH$, and let $\nu_n$ be the
probability law with CDF $G_n$. The mean of this envelope law is
$e_n=\int_\R(H-G_n)$. The family is fixed throughout; constants
are permitted to depend on it. Write $\phi,\Phi$ for the standard
Gaussian density and CDF.

\begin{theorem}\label{log:thm:main}
With $c=(\kappa_+-\kappa_-)/(3v\sqrt{2\pi e})$, there is a finite
constant $K$ such that for all integers $n\ge1$,
\begin{equation}\label{log:eq:main}
 \boxed{W_1\bigl(\nu_n,N(c\ln n,vn)\bigr)\le K,
 \qquad |e_n-c\ln n|\le K.}
\end{equation}
No lattice or aperiodicity assumption is needed.
\end{theorem}

\begin{corollary}\label{log:cor:cost}
For any finite stationary kernel with common action entropy $h$ and
common positive information variance $v$, in bits,
\[
 \boxed{C_n=nh+\frac{\kappa_+-\kappa_-}{3v\sqrt{2\pi e}}\ln n+O(1),}
\]
where $\kappa_i=\E(-\log_2P_i(Y_i)-h)^3$. The same expansion holds
for the envelope mean and for exact seeds of entropy within a bounded
additive amount of that mean.
\end{corollary}
\begin{proof}[Proof of the corollary from Theorem~\ref{log:thm:main}]
The centered transcript-information envelope is $G_n$, and
$L_n=nh+e_n$. Use the finite-tree sandwich~\eqref{eq:sandwich}.
\end{proof}


\subsection{Uniform Gaussian amplification}
For $t>0$ denote the $N(0,vt)$ density by
\[
 \gamma_t(x)=\frac{1}{\sqrt{2\pi vt}}e^{-x^2/(2vt)}.
\]
Let $P$ be convolution with $N(0,v)$, so $P\gamma_t=\gamma_{t+1}$.
This is a comparison operator, not an additional available action.

\begin{lemma}[One-step Gaussian defect]\label{log:lem:defect}
There is $0\le D<\infty$ such that for all $t\ge1$ and $x\in\R$,
\begin{equation}\label{log:eq:defect}
 |U\gamma_t(x)-\gamma_{t+1}(x)|
 \le D t^{-3/2}\gamma_{8t}(x).
\end{equation}
\end{lemma}
\begin{proof}
Let $M=\max_i\max\{|y|:y\in\operatorname{supp}X_i\}$.
Centered Taylor expansion and equality of the second moments give
\[
 \left|A_i\gamma_t-\gamma_t-\frac v2\gamma_t''\right|(x)
 \le\frac{\E|X_i|^3}{6}
       \sup_{|a|\le M}|\gamma_t'''(x-a)|.
\]
Also $\partial_t\gamma_t=(v/2)\gamma_t''$, whence
\[
 \left|\gamma_{t+1}-\gamma_t-\frac v2\gamma_t''\right|(x)
 \le\frac12\sup_{t\le s\le t+1}|\partial_s^2\gamma_s(x)|.
\]
For finite constants $D_3,D_4$, uniformly in $t\ge1$,
\begin{align}
 \sup_{|a|\le M}|\gamma_t'''(x-a)|
    &\le D_3t^{-3/2}\gamma_{8t}(x),\label{log:eq:window3}\\
 \sup_{t\le s\le t+1}|\partial_s^2\gamma_s(x)|
    &\le D_4t^{-2}\gamma_{8t}(x).\label{log:eq:window4}
\end{align}
Here is a direct justification of uniformity in $x$. With
$z=x/\sqrt{vt}$, $b=a/\sqrt{vt}$, one has $|b|\le M/\sqrt v$
and $(z-b)^2\ge z^2/2-b^2$. Gaussian spatial derivatives are a
polynomial times their Gaussian density. Thus the ratio in
\eqref{log:eq:window3}, after removing $t^{-3/2}$, is bounded by a
constant times $(1+|z|)^3e^{-3z^2/16}$. For~\eqref{log:eq:window4},
use $\partial_s^2\gamma_s=(v^2/4)\gamma_s''''$ and
$1\le s/t\le2$; the analogous ratio is bounded by a constant times
$(1+|z|)^4e^{-3z^2/16}$. Both suprema are finite.
Finally $t^{-2}\le t^{-3/2}$ and
$|\max_i a_i-b|\le\max_i|a_i-b|$, proving~\eqref{log:eq:defect}.
\end{proof}

\begin{lemma}[Bounded integrated amplification]\label{log:lem:amp}
Choose $A\ge1$ with $3D/\sqrt A\le1/2$, where $D$ is from
Lemma~\ref{log:lem:defect}. Then for every $t\ge A$, every $m\in\R$
and every integer $n\ge0$,
\begin{equation}\label{log:eq:amp}
 \boxed{\int_\R U^n\bigl[\gamma_t(\cdot-m)\bigr](x)\,dx\le2.}
\end{equation}
For example $A=\max\{1,36D^2\}$ is admissible.
\end{lemma}
\begin{proof}
Translation invariance reduces the proof to $m=0$. Positivity,
subadditivity,~\eqref{log:eq:lip}, and $P^r\gamma_t=\gamma_{t+r}$ give
the pointwise telescoping bound
\begin{equation}\label{log:eq:telescope}
 U^n\gamma_t\le\gamma_{t+n}
 +D\sum_{r=0}^{n-1}(t+r)^{-3/2}
                   U^{n-1-r}\gamma_{8(t+r)}.
\end{equation}
For clarity, obtain this by comparing $U^n\gamma_t$ with
$U^{n-1}\gamma_{t+1}$, then $U^{n-2}\gamma_{t+2}$, and so forth;
each difference is bounded by the appropriate iterate of $U$ applied
to $|U\gamma_{t+r}-\gamma_{t+r+1}|$.

Use strong induction on $n$, simultaneously for every real $t\ge A$.
The assertion at $n=0$ is the Gaussian integral, equal to one.
In~\eqref{log:eq:telescope}, every remaining horizon is less than $n$
and every Gaussian time $8(t+r)$ is at least $A$. The induction
hypothesis therefore bounds its integral by two. Since
\[
 \sum_{r=0}^\infty(t+r)^{-3/2}
 \le t^{-3/2}+\int_0^\infty(t+s)^{-3/2}\,ds
 \le\frac3{\sqrt t},
\]
the integrated right side of~\eqref{log:eq:telescope} is at most
$1+6D/\sqrt t\le2$. This closes the induction. There is no assumption
that the supremum over all horizons is already finite.
\end{proof}

\paragraph{A general stability consequence.}
If $|f-g|\le b\gamma_t(\cdot-m)$ with $b\ge0$ and $t\ge A$, then
\begin{equation}\label{log:eq:stab}
 \norm{T^nf-T^ng}\le2b\qquad(n\ge0).
\end{equation}
This follows from~\eqref{log:eq:lip} and Lemma~\ref{log:lem:amp}.
The same holds for sums of such majorants, with $2\sum b$ on the
right. A further consequence, for $f\ge0$ integrable, is
$\int U^n(\gamma_t*f)\le2\norm f$ when $t\ge A$:
take the supremum over policies inside the nonnegative integral
against $f$, and apply Tonelli and~\eqref{log:eq:amp}.


\subsection{The resonant corrector}
Define
\begin{align}
 f(z)&=\min_i\left[-\frac{\kappa_i}{6v^{3/2}}
                      (z^2-1)\phi(z)\right],\label{cor:eq:forcing}\\
 a_0&=\int_\R f(z)\,dz=-\frac{\kappa_+-\kappa_-}{3v^{3/2}}\phi(1)
       =-\frac{c}{\sqrt v},\label{cor:eq:a}\\
 g(z)&=2(a_0\phi(z)-f(z)),\qquad
 J(z)=\int_{-\infty}^z g(y)\,dy,\label{cor:eq:J}\\
 B(z)&=\phi(z)\int_0^z\frac{J(y)}{\phi(y)}\,dy.\label{cor:eq:B}
\end{align}
Indeed $\int(z^2-1)\phi=0$ and
$\int|z^2-1|\phi=4\phi(1)$, proving~\eqref{cor:eq:a}.
Since $\int g=0$, formula~\eqref{cor:eq:B} decays at both ends and solves
\begin{equation}\label{cor:eq:ODE}
 B''+zB'+B=2(a_0\phi-f).
\end{equation}

\paragraph{Regularity and tails, explicitly.}
The function $f$ is continuous, smooth off $\{-1,1\}$, and has only
derivative jumps there. Thus $B\in C^2$, with locally absolutely
continuous second derivative. For $z<-1$ and $z>1$, respectively,
\[
 \frac{J(z)}{\phi(z)}=
 2a_0\frac{\Phi(z)}{\phi(z)}-\frac{\kappa_+z}{3v^{3/2}},
 \qquad
 \frac{J(z)}{\phi(z)}=
 -2a_0\frac{1-\Phi(z)}{\phi(z)}-\frac{\kappa_+z}{3v^{3/2}}.
\]
The elementary Gaussian tail inequality
$1-\Phi(z)\le\phi(z)/z$ for $z>0$, and its reflected version,
show that the integral in~\eqref{cor:eq:B} is $O(1+z^2)$.
The equation $B'+zB=J$ and~\eqref{cor:eq:ODE}, followed by its
almost-everywhere derivative, then show that $B,B',B'',B'''$ are
bounded by a polynomial times $\phi$ on each tail. All their
polynomially weighted absolute integrals are finite. In particular
$B\in W^{3,1}(\R)$. No fourth classical derivative of $B$ is required.

For $t\ge1$, put
\begin{equation}\label{cor:eq:profile}
 z=z_t(x)=\frac{x-c\ln t}{\sqrt{vt}},\qquad
 p_t(x)=\Phi(z),\quad q_t(x)=t^{-1/2}B(z),\quad u_t=p_t+q_t.
\end{equation}
The function $u_t$ need not be a CDF; it is used only as a bounded
comparison profile. Its limiting values are zero and one. Direct
differentiation using~\eqref{cor:eq:ODE} gives
\begin{equation}\label{cor:eq:PDE}
 \partial_tu_t-\frac v2\partial_x^2u_t
 =t^{-3/2}f(z)-\frac{c}{\sqrt v}\,t^{-2}B'(z).
\end{equation}
The translation $c\ln t$ absorbs the resonant Gaussian mode exactly.


\subsection{Pointwise discrete residual}\label{sec:residual}
The following estimate closes the analytic input needed for bounded
propagation. All constants below depend only on the fixed increment family.
Put
\[
 E_0(z)=(1+z^2)^4e^{-z^2/8}.
\]
\begin{lemma}\label{log:lem:residual}
There is $K_0<\infty$ such that for every real $t\ge1$ and $x\in\R$,
\begin{equation}\label{log:eq:input}
 |u_{t+1}(x)-Tu_t(x)|\le K_0t^{-2}E_0(z_t(x)).
\end{equation}
\end{lemma}
\begin{proof}
Write $p_t(x)=\Phi(z_t(x))$ and $q_t(x)=t^{-1/2}B(z_t(x))$.
We first specify the uniform Gaussian bounds used in Taylor's formula.
The corrector construction gives
\[
 |B^{(j)}(z)|\le C(1+|z|)^{j+2}\phi(z),\qquad 0\le j\le3,
\]
where the last derivative is interpreted almost everywhere. For $j=0$
this follows from the integral representation and the Gaussian tail bound;
$B'=J-zB$, $B''=2(a_0\phi-f)-zB'-B$, and the almost-everywhere derivative
of the last identity give $j=1,2,3$. On the compact interval $[-1,1]$
all these bounds follow by boundedness, enlarging $C$ if necessary.

Let $M$ bound every increment. For a spatial shift $|y|\le M$, the
standardized argument changes from $z=z_t(x)$ by at most $M/\sqrt v$.
For $t\le s\le t+1$ the time-changed argument is
\[
 z_s(x)=\sqrt{t/s}\,z+
       \frac{c\ln(t/s)}{\sqrt{vs}}.
\]
Here $\sqrt{t/s}\in[1/\sqrt2,1]$, and the additive term is uniformly
bounded. The inequality $(r z+d)^2\ge r^2z^2/2-d^2$, valid for all real
$z,d$, shows that any polynomial of degree at most eight times
$\phi(rz+d)$ is bounded by a constant times $E_0(z)$ in these windows.
For spatial shifts use $r=1$; for time changes use $r\ge1/\sqrt2$.
Thus the constants in the following derivative estimates are uniform
in $t,x$ and throughout the indicated Taylor segments:
\begin{align*}
 \sup_{|y|\le M}|p_t^{(4)}(x-y)|&\le Ct^{-2}E_0(z),\\
 \operatorname*{ess\,sup}_{|y|\le M}|q_t^{(3)}(x-y)|
      &\le Ct^{-2}E_0(z),\\
 \sup_{s\in[t,t+1]}|\partial_s^2u_s(x)|&\le Ct^{-2}E_0(z).
\end{align*}
For the last line, differentiate
$z_t'=-z_t/(2t)-c/(\sqrt v\,t^{3/2})$ once more. The resulting terms
involve only $\phi,\phi',B,B',B''$, multiplied by polynomials of degree
at most two and by powers no larger than $t^{-2}$ for $t\ge1$.
The preceding window estimate applies to every term.

Centered spatial Taylor expansion, first to order three and then to
order two, gives
\begin{align*}
 A_i p_t&=p_t+\frac v2p_t''-\frac{\kappa_i}{6}p_t'''+r_{i,t},\\
 A_i q_t&=q_t+\frac v2q_t''+s_{i,t},
\end{align*}
with $|r_{i,t}(x)|+|s_{i,t}(x)|\le Ct^{-2}E_0(z)$.
Indeed the two integral remainders are bounded by
$\E|X_i|^4/24$ and $\E|X_i|^3/6$ times the respective derivative
suprema. The $C^2$ regularity and local absolute continuity of $B''$
justify Taylor's formula across the switching points $\pm1$.
Taking the minimum over $i$ is Lipschitz in the maximum norm, so
\[
 Tu_t=u_t+\frac v2u_t''+t^{-3/2}f(z)+O(t^{-2}E_0(z)).
\]
On the other hand, the corrector equation gives exactly
\[
 \partial_tu_t-\frac v2u_t''
   =t^{-3/2}f(z)-\frac c{\sqrt v}t^{-2}B'(z).
\]
Time Taylor expansion on $[t,t+1]$ and the time-derivative estimate above yield
\[
 u_{t+1}=u_t+\frac v2u_t''+t^{-3/2}f(z)+O(t^{-2}E_0(z)).
\]
Subtracting proves the claim.
\end{proof}

Let $D_E=64^4e^{-63/16}$. Maximizing
$(1+z^2)^4e^{-z^2/16}$ shows that it is at most $D_E$. Consequently,
with $b_0=K_0D_E\sqrt{16\pi v}$,
\begin{equation}\label{log:eq:gaussres}
 |u_{t+1}-Tu_t|(x)\le b_0t^{-3/2}\gamma_{8t}(x-c\ln t).
\end{equation}
The mass of this majorant is summable over $t=T,T+1,\ldots$.
The amplification lemma ensures that summability survives every remaining
Bellman iteration.

\subsection{Completion of the bounded-remainder proof}
Fix any real $T\ge\max\{A,2\}$ once and for all, independently of $n$.
The initial error has a Gaussian majorant
\begin{equation}\label{log:eq:initial}
 |H(x)-u_T(x)|\le b_1\gamma_{s_0}(x-c\ln T)
\end{equation}
for some finite $b_1$ and some $s_0\ge A$.
To verify this, choose $s_0=8T$. On each sufficiently distant tail,
$H-\Phi(z_T)$ is a Gaussian tail and $B(z_T)$ is bounded by a
polynomial times $\phi(z_T)$. Dividing by the broader Gaussian
$\gamma_{8T}(x-c\ln T)$ gives a bounded ratio tending to zero.
On the remaining compact interval the numerator is bounded and
the denominator has a positive minimum. The jump of $H$ at zero
does not affect this pointwise bound.

Using~\eqref{log:eq:lip} repeatedly, followed by~\eqref{log:eq:amp},
\eqref{log:eq:gaussres}, and~\eqref{log:eq:initial}, gives
\begin{align}
 \norm{G_n-u_{T+n}}
 &\le\int U^n|H-u_T|
    +\sum_{j=0}^{n-1}\int U^{n-1-j}|u_{T+j+1}-Tu_{T+j}|\notag\\
 &\le2b_1+2b_0\sum_{j=0}^{n-1}(T+j)^{-3/2}\notag\\
 &\le2b_1+\frac{6b_0}{\sqrt T}=:K_1.\label{log:eq:bounded}
\end{align}
This is uniform in $n$. Moreover
\[
 \norm{u_t-\Phi(z_t)}=\sqrt v\norm B.
\]
Finally the Gaussians with means $c\ln(n+T)$ and $c\ln n$,
and variances $v(n+T)$ and $vn$, can be coupled using the same
standard normal variable. Their $W_1$ distance is at most
\[
 |c|\ln(1+T/n)
 +\sqrt{2v/\pi}\,(\sqrt{n+T}-\sqrt n),
\]
which tends to zero and is bounded for $n\ge1$.
The one-dimensional formula $W_1(\mu,\nu)=\int|F_\mu-F_\nu|$
and~\eqref{log:eq:bounded} prove the first part of~\eqref{log:eq:main}.
The difference of means is bounded by $W_1$, proving the second.
This proves Theorem~\ref{log:thm:main}.\hfill$\square$



\subsection{An explicit information pair}
Take
\[
 P=(1/4,\underbrace{1/16,\ldots,1/16}_{12}),\qquad
 Q=(\underbrace{1/8,\ldots,1/8}_{6},
       \underbrace{1/32,\ldots,1/32}_{8}).
\]
For $P$, centered information is $-3/2$ with probability $1/4$ and
$1/2$ with probability $3/4$. For $Q$, it is $-1/2$ with probability
$3/4$ and $3/2$ with probability $1/4$. Direct calculation yields
\[
 h=7/2,\qquad v=3/4,\qquad \kappa_P=-3/4,\quad\kappa_Q=3/4.
\]
Therefore
\[
 \boxed{C_n=\frac72 n+\frac{2}{3\sqrt{2\pi e}}\ln n+O(1).}
\]
This is the optimum over all exact causal seeds, not the cost of a
restricted block architecture. The coefficient is
$0.161313816346\ldots$ bits per unit of $\ln n$.
The example's dyadic structure is not used in proving this asymptotic.
A separate dyadic refinement theorem in the repository additionally
identifies $C_n=L_n$ at every horizon.

\subsection{What bounded amplification resolves}
A summable local residual alone would be insufficient if its propagation
could grow with the remaining horizon. Lemma~\ref{log:lem:amp} controls
precisely that growth on the Gaussian majorants produced by the corrector.
A fixed starting time then suffices. No smoothness of the actual envelope
$G_n$, nonlattice hypothesis, or prior logarithmic bound on its width is
needed in this proof. In particular, when the third moments all agree,
$c=0$ and the same argument proves bounded excess directly.

\section{Formalization, reproducibility, and scope}\label{sec:formal}
The formal development is available at
\begin{center}\url{https://github.com/jendamikulik/finite-Shannon-theorem}.\end{center}
The revision used for this paper is
\begin{center}\texttt{5741ee193bd5ff7f4addd4c09b2d65edfe238c20}.\end{center}
It pins Lean \texttt{4.35.0-rc3} and Mathlib revision
\path{7d6757bc18680044fc0ff6efc8ee20494b408556}.
The import root is \path{CausalSeed/Main.lean}. The repository supplies the
commands
\begin{verbatim}
lake build
lake env lean scripts/AxiomAudit.lean
\end{verbatim}
for compilation and enumeration of theorem dependencies.
The recorded \path{AXIOMS.txt} lists 2528 project theorems; all recorded
dependencies are among \texttt{propext}, \texttt{Classical.choice}, and
\texttt{Quot.sound}, and seven theorems use none. It reports no nonstandard
axioms and no project axiom declarations. The audit includes
\path{Kernel.cost_bounded} and \path{LogGauss.bounded_L1}.
These are the recorded results for the cited revision, not a report of a
new build performed while typesetting this manuscript.

\subsection{Statement-to-source map}
All names below lie under \path{CausalSpectrum}. Paths are relative to
\path{CausalSeed/}. Names are provided instead of unstable line numbers.
\begingroup\small
\begin{longtable}{p{0.26\textwidth}p{0.68\textwidth}}
\toprule
Mathematical statement & Module and principal declaration \\
\midrule\endhead
Finite Shannon bound & \path{Entropy.lean}: \path{one_seed_shannon}, \path{causal_realization} \\[3pt]
Information converse & \path{Penalty.lean}: \path{entropy_ge_envelope}; \path{Entropy.lean}: \path{seedCDF_le_Fstar} \\[3pt]
CDF/transport identity & \path{GapW1.lean}: \path{wasserstein_gap_W1}; \path{Wasserstein.lean}: \path{OT.W1_eq_cdfDist} \\[3pt]
Projective identity & \path{EntropyLimit.lean}: \path{Completion.projective_entropy_identity}, \path{Completion.entropy_attained} \\[3pt]
Infinite sandwich & \path{InfiniteEnvelope.lean}: \path{Completion.infinite_envelope_sandwich} \\[3pt]
Atomic capacity & \path{AtomicCapacity.lean}: \path{Completion.causal_atomic_capacity_attained}, \path{Completion.countable_seed_iff_atomQ_one} \\[3pt]
Capacity from the envelope & \path{AtomicSpectrum.lean}: \path{atomQ_eq_lim_FInf} \\[3pt]
Simultaneous R\'enyi bound & \path{RenyiLimitMeasure.lean}: \path{RenyiInf.one_seed_all_orders_meas} \\[3pt]
Zero branches & \path{ZeroClosure.lean}; \path{ZeroBranchTransfer.lean}; \path{ZeroBranchOrdinary.lean} \\[3pt]
Square-root asymptotic & \path{SqrtGeneral.lean}: \path{Kernel.cost_tendsto}; \path{SqrtMain.lean}: \path{Kernel.cost_rate} \\[3pt]
Universality & \path{SqrtW1.lean}: \path{Kernel.universality_W1} \\[3pt]
Corrector and residual & \path{LogCorrector.lean}: \path{LogGauss.cB_ode}; \path{LogResidual.lean}: \path{LogGauss.residual_pointwise} \\[3pt]
Gaussian amplification & \path{LogBounded.lean}: \path{LogGauss.amp_bounded} \\[3pt]
Bounded logarithmic remainder & \path{LogBounded.lean}: \path{LogGauss.bounded_L1}, \path{LogGauss.bounded_mean}, \path{Kernel.cost_bounded} \\[3pt]
Explicit dyadic pair & \path{LogBounded.lean}: \path{Kernel.LogPair.cost_bounded} \\
\bottomrule
\end{longtable}\endgroup

\subsection{Definitions and the verification boundary}
The finite objects \path{Tree}, \path{Policy}, \path{Strategy},
\path{run}, and \path{leafMass} are defined in \path{Spectrum.lean}.
The functions \path{Fstar} and \path{envelopeMean} in \path{Entropy.lean}
represent the information envelope and its mean. The infinite model uses
\path{ITree}, \path{Completion.CountableSeed}, and its predicates
\path{ExactTo} and \path{Exact}. Ranking and finite restriction are
connected to arbitrary countable seed alphabets by the bridge modules.
The infinite implementation uses extended entropy in nats;
\path{Kernel.costBits} and the envelope bridges convert to bits.

An axiom audit checks logical dependencies of the formal statements.
Reading the definitions and their hypotheses is still necessary to assess
whether those statements express the intended simulation problem. In
particular, the minimum over transcript CDFs is taken pointwise; it is
not assumed to be the law of one policy. The fixed-kernel constants in
Section~\ref{sec:log} are existential and need not be uniform as the
variance or the kernel varies. Formal verification does not establish
historical priority or numerical practicality.

\subsection{Open questions}
The bounded logarithmic remainder need not be treated as a convergent
constant. This paper proves neither convergence of $C_n-nh-c\ln n$ nor
an expansion for its successive differences. Even if the envelope mean
has a limiting constant, the bounded entropy gap must still be controlled
to transfer such a limit to $C_n$.

Other questions in the accompanying development include the optimal
universal additive constants and the exact supremum $G_*$ of the causal
minus ordinary Shannon optimum. Its recorded bounds are
$2/9\le G_*\le c_\star$; the endpoint behavior of exact penalties at
$G_*$ is separate from the asymptotic theorem here. For unequal action
entropies, an $O(\ln n)$ comparison with the maximal-entropy subfamily
does not transfer a logarithmic coefficient. We therefore impose common
entropy throughout the sharp logarithmic theorem.

\paragraph{Acknowledgment of assistance.}
AI systems, including ChatGPT and Grok, were used during the development
of mathematical arguments, code, formalization, and exposition. The cited
source revision and the explicit statements above delimit the mathematical
claims; the use of these systems is not evidence of their correctness.

\begin{thebibliography}{99}
\raggedright
\bibitem{compton}
S. Compton, D. Katz, B. Qi, K. Greenewald, and M. Kocaoglu.
Minimum-Entropy Coupling Approximation Guarantees Beyond the Majorization
Barrier. In \emph{Proceedings of AISTATS}, PMLR \textbf{206},
10445--10469, 2023.
\url{https://proceedings.mlr.press/v206/compton23a.html}.

\bibitem{shkel}
Y. Y. Shkel and A. K. Yadav.
\emph{Information Spectrum Converse for Minimum Entropy Couplings and
Functional Representations}. 2023 IEEE International Symposium on
Information Theory; preprint arXiv:2305.05745.
\url{https://arxiv.org/abs/2305.05745}.

\bibitem{li}
C. T. Li. Efficient Approximate Minimum Entropy Coupling of Multiple
Probability Distributions. \emph{IEEE Transactions on Information Theory}
\textbf{67}(8), 5259--5268, 2021.
\href{https://doi.org/10.1109/TIT.2021.3076986}{doi:10.1109/TIT.2021.3076986}.

\bibitem{ptas}
S. Compton. \emph{Efficient $\varepsilon$-approximate minimum-entropy
couplings}. Preprint arXiv:2509.19598, 2025.
\url{https://arxiv.org/abs/2509.19598}.

\bibitem{peng}
S. Peng. \emph{A New Central Limit Theorem under Sublinear Expectations}.
Preprint arXiv:0803.2656, 2008.
\url{https://arxiv.org/abs/0803.2656}.

\bibitem{pz}
S. Peng and Q. Zhou. \emph{A hypothesis-testing perspective on the
$G$-normal distribution theory}. Preprint arXiv:1909.03530, 2019.
\url{https://arxiv.org/abs/1909.03530}.

\bibitem{fang}
X. Fang, S. Peng, Q.-M. Shao, and Y. Song.
\emph{Limit theorems with rate of convergence under sublinear expectations}.
Preprint arXiv:1711.10649, 2017, revised 2018.
\url{https://arxiv.org/abs/1711.10649}.

\bibitem{repo}
J. Mikul\'ik. \emph{Finite Shannon theorem}. Lean source repository,
revision \texttt{5741ee193bd5ff7f4addd4c09b2d65edfe238c20}, 12 October 2026.
\url{https://github.com/jendamikulik/finite-Shannon-theorem/tree/5741ee193bd5ff7f4addd4c09b2d65edfe238c20}.
\end{thebibliography}
\end{document}
